\begin{twocolentry}{}
\href{https://github.com/faketut/MiniKV}{\textbf{MiniKV}}
\end{twocolentry}
\begin{onecolentry}
\begin{highlights}
\item 实现精简版 LSM-tree 存储：追加写 WAL、内存 memtable 刷盘为 SSTable，以及通过 compaction 合并重叠的有序 run，以控制读放大。
\item 增加 SSTable block 索引与清晰的恢复路径（重放 WAL、丢弃不完整尾部），确保写入过程中即便被 kill -9 也不会破坏已提交的键。
\item 调优 flush/compaction 触发条件，并在读路径复用缓冲区；相较“每次写入即刷盘”的朴素设计，提升并发写入与范围扫描下的吞吐。
\end{highlights}
\end{onecolentry}

\begin{twocolentry}{}
\href{https://github.com/faketut/MiniKV}{\textbf{MiniKV}}
\end{twocolentry}
\begin{onecolentry}
\begin{highlights}
\item 实现轻量级 LSM-tree 键值存储：追加写 WAL、内存 memtable 刷盘生成 SSTable，并通过 compaction 合并重叠的有序 run，以限制读放大。
\item 增加 SSTable 分块索引与清晰的恢复流程（回放 WAL、丢弃不完整尾部），确保写入过程中遭遇 kill -9 也不会破坏已提交的键。
\item 调整刷盘/压缩触发阈值，并在读路径复用缓冲区，相比“每次写入即刷盘”的朴素实现，在并发写入与范围扫描下提升吞吐。
\end{highlights}
\end{onecolentry}
